% Source: https://artofproblemsolving.com/community/c3h2143358_a_problem
% Stable identity: aops-topic-2143358-q01
Three mutually tangent circles of diameter $1$ have centres forming an equilateral triangle of side $1$, whose height is $\dfrac{\sqrt3}{2}$. So consecutive rows of centres are $\dfrac{\sqrt3}{2}$ apart.

\begin{asy}
usepackage("amsmath");
size(7cm);
defaultpen(fontsize(10pt));
real h=sqrt(3)/2;
for(int k=0;k<4;++k){ real y=0.5+k*h; real x0=(k%2==0)?0.5:1.0;
  for(int i=0;i<6;++i){ real x=x0+i; if(x<5.6) draw(circle((x,y),0.5),gray(0.55)); } }
draw((0,0)--(6.2,0),linewidth(1.3)); draw((0,0)--(0,3.9),linewidth(1.3));
pair P1=(1.5,0.5), P2=(2.5,0.5), P3=(2.0,0.5+h);
draw(P1--P2--P3--cycle,red+linewidth(1)); draw(P3--(2.0,0.5),red+dashed);
dot(P1,red); dot(P2,red); dot(P3,red);
label("$1$",(2.0,0.45),S,red+fontsize(9pt),Fill(white));
label("$\dfrac{\sqrt3}{2}$",(2.5,0.5+h/2),E,red+fontsize(9pt),Fill(white));
label("\mbox{}",(min(currentpicture,user=true)+max(currentpicture,user=true))/2,black); // keeps the text after the figure black when LaTeX typesets the labels inline
\end{asy}

\textbf{Number of rows.} The centres must lie at least $\dfrac{1}{2}$ from each side of the square, so the centre lines of the rows lie in a band of width $99$. With $r$ rows we need $(r-1)\cdot\dfrac{\sqrt3}{2}\le99$, that is,
\[r-1\le\dfrac{198}{\sqrt3}=66\sqrt3\approx114.3,\]
so at most $r=115$ rows fit, and $115$ rows do fit.

\textbf{Circles per row.} A row holds at most $100$ circles, with centres at $\dfrac{1}{2},\dfrac{3}{2},\dots,\dfrac{199}{2}$ from the left side. A row next to a full row has its centres shifted by $\dfrac{1}{2}$, so they are at whole numbers, and only $1,2,\dots,99$ are allowed: that row holds $99$ circles. In general, two neighbouring rows cannot both hold $100$ circles, so any two neighbouring rows hold at most $199$ circles.

\textbf{Best arrangement.} With $115$ rows, alternate rows of $100$ and $99$, starting and ending with a row of $100$: $58$ rows of $100$ and $57$ rows of $99$, in total
\[58\cdot100+57\cdot99=5800+5643=11443\]
circles. This is the most possible: pairing the first $114$ rows into $57$ neighbouring pairs gives at most $57\cdot199+100=11443$. (Fewer rows give fewer circles, since each row holds at most $100$ and every two neighbouring rows at most $199$.)

So the number of additional circles is $11443-10000=\boxed{1443}$.
